\documentclass[12pt,a4paper]{article}
\usepackage{geometry}
\usepackage{microtype}
\geometry{margin=2.8cm}
\usepackage{parskip}

\begin{document}
\thispagestyle{empty}

\noindent
Professor Salvatore Greco\\
Editor-in-Chief\\
\textit{Decisions in Economics and Finance}\\
Department of Economics and Business\\
University of Catania\\
Corso Italia 55, 95129 Catania, Italy

\vspace{1cm}

\noindent Dear Professor Greco,

\vspace{0.4cm}

Please find enclosed our manuscript entitled \textbf{``Loss Aversion, Endogenous
Reference Points, and Boom-Bust Asymmetry in Financial Wealth Dynamics''} for
consideration for publication in \textit{Decisions in Economics and Finance}.

\textbf{Summary of contribution.} The paper derives boom-bust asymmetry in
financial wealth dynamics as an equilibrium theorem from a minimal
CPT-like structure. We retain the single empirically robust element of
Cumulative Prospect Theory --- the kink in the value function at a reference
point --- and embed it in a Stackelberg market with a loss-averse market maker,
a rational speculator, and a shared Schumpeterian wealth benchmark. Two
structural features are jointly necessary: the zero-sum trading game guarantees
that the two agents always sit on opposite sides of the kink simultaneously,
and the Schumpeterian trend pins volatility persistence to the economy's
long-run growth rate. The paper derives a structural Threshold-GARCH recursion
for the conditional variance of detrended wealth deviations with a separation
property: the persistence parameter is determined by the growth rate alone and
the asymmetry ratio by dealer loss aversion alone. All structural mappings are
derived analytically and verified symbolically. An out-of-sample test
using the 2008 Global Financial Crisis as a natural experiment confirms
all three structural predictions: the model estimated on pre-GFC data
alone correctly predicts the asymmetric compression of post-GFC variance,
stagnation-period persistence, and the sign reversal of the variance
ratio, with zero post-GFC parameters.

\textbf{Fit with the journal.} The paper sits at the intersection of
mathematical finance, behavioural economics, and market microstructure ---
the core scope of \textit{Decisions in Economics and Finance}. It is
theory-driven, derives closed-form results, and connects to an empirically
testable cross-sectional identification strategy.

\textbf{Declarations.} This manuscript has not been published previously and
is not under consideration for publication elsewhere. The author declares no
competing interests.

\vspace{0.6cm}

\noindent Yours sincerely,

\vspace{0.8cm}

\noindent Anurag Srivastava\\
Riskcare Ltd., London, United Kingdom\\
ORCID: 0000-0002-6477-4430

\end{document}
